\begin{table}[H]
\centering
\renewcommand{\arraystretch}{1.3}
\begin{tabular}{|p{4cm}|p{5cm}|p{5cm}|}
\hline
\textbf{Criterio} & \textbf{Operaciones Sencillas} & \textbf{Operaciones Complejas} \\ \hline
\textbf{Entorno de operación} & 
Zonas rurales o semi-rurales, sin multitudes ni infraestructura crítica. & 
Áreas urbanas densamente pobladas, eventos masivos o cercanía a aeropuertos. \\ \hline

\textbf{Condiciones de vuelo} & 
VLOS (Visual Line of Sight), un único UAS, condiciones meteorológicas favorables (VMC). & 
BVLOS o EVLOS, múltiples UAS, condiciones meteorológicas limitantes que requieren mitigación. \\ \hline

\textbf{Altura y espacio aéreo} & 
Inferiores a 400 ft AGL, fuera de espacio aéreo controlado. & 
Operaciones en espacio aéreo controlado o dentro de áreas con restricciones especiales. \\ \hline

\textbf{Personal mínimo requerido} & 
Jefe de Pilotos (puede actuar como piloto) + 1 Observador. & 
Jefe de Pilotos + Pilotos adicionales + Observadores según complejidad. \\ \hline

\textbf{Participación del Gerente SMS} & 
No está presente en campo. Revisa registros y checklists posterior a la operación. & 
Supervisa análisis de riesgos, valida mitigaciones y puede estar presente en campo si se estima necesario. \\ \hline

\textbf{Nivel de documentación} & 
Checklists de pre-vuelo, vuelo y post-vuelo. Registro en bitácora. & 
Análisis de riesgos formal, plan de mitigación, autorización especial y registros adicionales. \\ \hline
\end{tabular}
\caption{Cuadro comparativo de clasificación de operaciones UAS en Skyfilms.}
\end{table}
